
\subsection{Checked Lean steps}

These entries are proofs of the specified steps. They supplement
the manuscript-to-Lean statement map. In particular, an encoded
target proposition with an open proof has no \verb|\leanok| tag.

\begin{lemma}[Finite norms and Lipschitz circle Fourier series]\label{checked-norm-foundations}
\lean{NilBohr.intervalGowersForm_norm_le,NilBohr.normalizedGowers_triangle,NilBohr.normalizedGowersAdjacentMonotonicity,NilBohr.NormalizedGowersLaws.checked,NilBohr.progressionGowers_one_eq_norm_mean,NilBohr.progressionGowers_le_of_le,NilBohr.phaseNorm_le_progressionGowers,NilBohr.normalizedGowers_dilation,NilBohr.kmttNormalizedDilation,NilBohr.bernsteinCircleLipschitzFourier,NilBohr.fourierCoeff_tail_eta_le,NilBohr.norm_fourierSeries_truncation_eta_le}\leanok
The interval cube norm obeys the triangle inequality for every positive
order. On every nonempty finite arithmetic progression, normalized
order comparison and $u^s\ll_s U^s$ have uniform constants. In the
change of variable $x=qy+b$ the normalized dilation bound loses at
most an order-dependent constant times $q^{(s+1)/2^s}$. Lipschitz
circle cutoffs have absolute Fourier mass bounded by a universal
constant times their supremum norm plus Lipschitz seminorm.
\end{lemma}
\begin{proof}
\proves{checked-norm-foundations}\leanok
Coordinate Cauchy--Schwarz bounds the mixed cube form. A separate
direction-sum Cauchy--Schwarz and indicator cube counts give the
normalized order comparison. Exact affine cube transport handles
progressions, and a finite residue Fourier expansion handles dilation.
Parseval applied to translated differences, followed by a limit on
each finite frequency set, bounds the weighted Fourier square sum.
Cauchy--Schwarz against the reciprocal frequencies gives absolute
summability. The reciprocal-square tail is at most $4/R$ outside
$|h|\le R$, giving an $O(\operatorname{Lip}(F)\eta)$ absolute
Fourier tail at the manuscript's cutoff $R=\eta^{-2}$.
\end{proof}

\begin{lemma}[Fixed local factors from the inverse theorem]\label{checked-weak-regularity}
\lean{NilBohr.KMTTLocalInverse.toLocalWeakRegularity,NilBohr.exists_localInverse_fixed_coordinate,NilBohr.fixedStep_progression_partition}\leanok
The local inverse estimate supplies the weak-regularity interface
used in Proposition~3.1, including the partition, progression length
and small-scale alternative.
\end{lemma}
\begin{proof}
\proves{checked-weak-regularity}\uses{external-inputs}\leanok
Choose a fixed power of $q$ divisible by every eligible inverse step.
Partition its residue classes into nearly equal blocks at the
coordinate-dependent resolutions. A telescoping discrepancy is witnessed
by a normalized residual in one coordinate. Reanchor at $P_1$ when
the residual is in the base coordinate. Each atom has zero residual
mean, so a function with the inverse Lipschitz bound has correlation
less than the lower bound forced by the inverse theorem.
\end{proof}

\begin{lemma}[Equivalent prime number theorem limits]\label{checked-prime-normalization}
\lean{NilBohr.PrimeNumberTheorem.of_theta,NilBohr.PrimeNumberTheorem.of_psi,NilBohr.PrimeNumberTheorem.of_count,NilBohr.PrimeNumberTheorem.to_psi,NilBohr.logarithmicIntegral_normalization_tendsto,NilBohr.PrimeNumberTheoremInProgressions.of_logarithmicIntegral}\leanok
The theta, psi and normalized ordinary prime-count asymptotics give
equivalent premises for the conversion to ordinary prime averages.
The fixed-modulus prime count can use either $\operatorname{Li}(N)$
or $N/\log N$ as its scale.
\end{lemma}
\begin{proof}
\proves{checked-prime-normalization}\leanok
The contribution of higher prime powers is $o(N)$. For any fixed
$0<a<1$, the primes at most $N^a$ contribute $o(1)$ to the normalized
count, and on the tail $a\log N\le\log p\le\log N$. Taking $a$ close
to one proves the equivalence with the ordinary count limit.
The same split for the integral of $1/\log t$ on $[2,N]$ proves the
equivalence of the fixed-modulus scales.
\end{proof}

\begin{lemma}[The Heisenberg return window]\label{checked-heisenberg-window}
\lean{NilBohr.heisenbergWindow,NilBohr.isOpen_heisenbergWindow,NilBohr.heisenbergWindow_measure,NilBohr.heisenbergRotation_mem_window_iff,NilBohr.heisenbergWindow_basicNilsequence,NilBohr.heisenbergWindowNilsequence_re_pos_iff,NilBohr.heisenbergWindowNilsequence_prime_re_pos_iff}\leanok
The window $0<y<1$, $z\bmod 1\in J$ defines an open subset of the
right-coset quotient and has Haar mass $|J|$. For admissible $\beta$
and each nonzero integer $n$, the orbit at time $n$ belongs to this
window exactly when $Z(n)\bmod 1\in J$. A basic polynomial nilsequence
realizes the window by the open right half-plane in $\C$.
\end{lemma}
\begin{proof}
\proves{checked-heisenberg-window}\uses{selector}\leanok
Pull the window back to the lattice translates of the open strip.
In the unit Haar fundamental domain its volume is $|J|$ by Fubini.
The signed group powers give the orbit formula. A nonzero integer
multiple of $\beta$ avoids the seam by irrationality. The exact
orbit membership criterion includes the seam condition at time zero.
The distance to the closed complement is a continuous function on
the quotient, and the signed powers are polynomial for the rational
lower central filtration. Complexifying the distance supplies the
required continuous test.
\end{proof}

\begin{lemma}[Prime distribution from the model and sieve]\label{checked-prime-distribution}
\lean{NilBohr.tendsto_cramerWeight_mean,NilBohr.ambient_prime_polynomial_test_error_tendsto,NilBohr.primeNumberTheorem_of_prime_model,NilBohr.primeNumberTheoremInProgressions_of_kmtt_and_sieve,NilBohr.tendsto_cramer_linear_mean,NilBohr.vinogradov_of_prime_model}\leanok
Banks's zero exclusion, the ambient Gowers model, Mertens' interval
bound and the fundamental sieve lemma supply the ordinary prime
number theorem, its fixed reduced progression form, and irrational
linear equidistribution over primes against continuous tests.
\end{lemma}
\begin{proof}
\proves{checked-prime-distribution}\uses{external-inputs,checked-norm-foundations,checked-prime-normalization}\leanok
Average the sieve bounds over a complete period to get the exact
Euler center. On the actual interval or the fixed reduced residue
class each divisor count has error at most one. The order-two
Gowers error transfers the means to Mangoldt, including the fixed
residue projection. For an irrational frequency, use reduced
Dirichlet approximants with denominator at most $\lfloor\sqrt N\rfloor$.
Compare the complex sieve midpoint with its normalized rational
Fourier coefficient over a complete period. Off the multiples of
the denominator the geometric sums are short. On the multiples,
the residual phase differs from its block mean by a controlled
error. The periodic coefficient is bounded by $\sqrt{2/q}$ and the
approximate denominator $q$ tends to infinity at all large scales.
Transfer the linear tests through the model and apply trigonometric
density after conversion to ordinary primes.
\end{proof}

\begin{lemma}[Modulus budgets in Appendix A]\label{checked-iteration-moduli}
\lean{NilBohr.log_iteration_modulus_le_density_power,NilBohr.eventually_density_iteration_moduli_small,NilBohr.eventually_small_modulus_absorbed_by_logarithmic_saving,NilBohr.eventually_large_admissible_modulus_at_polylog_density}\leanok
Given the Section~3 iteration budgets, $\log q_i\ll\delta^{-D_0}$
for a fixed $D_0$. If $\delta\ge(\log\log N)^{-c}$ and $cD_0<1$,
then $q_i\le(\log N_i)^\epsilon$ for every fixed $\epsilon>0$,
uniformly over stages. Every fixed power loss in $q_i$ can be
absorbed by any positive slack in a logarithmic saving. At the
polylogarithmic threshold the admissible numerical parameters allow
$q>(\log N)^A$ for each fixed $A\ge0$.
\end{lemma}
\begin{proof}
\proves{checked-iteration-moduli}\leanok
Bound the stage index using the geometric density increase and
use $\log N_i\ge(\log N)^{3/4}$. For the admissibility example
choose $w=\lceil K(\log N)^{\kappa E}\rceil$ with $\kappa E<c_{\rm rel}$
and take the largest power of $\mathcal L(w)$ with logarithm at
most $w^2$. Its logarithm exceeds $w^2/2$, and the prime support
and divisibility requirements follow from those of $\mathcal L(w)$.
\end{proof}

\begin{lemma}[Rescaling and the vertical chart]\label{checked-rescaling}
\lean{NilBohr.eval_rescalePolynomial,NilBohr.eval_rescalePolynomial_mul,NilBohr.HasCoefficients.rescale,NilBohr.heisenberg_chart_right_lattice}
\leanok
For a polynomial with zero constant coefficient, rescaling has the
prescribed coefficients and commutes with successive affine lifts.
The Heisenberg coordinate change under the right lattice is integral
in the vertical phase and preserves the horizontal fractional part.
\end{lemma}
\begin{proof}
\proves{checked-rescaling}\leanok
The polynomial calculation is exact in $\Z[y]$ and preserves degree
for nonzero scaling. The chart statement follows by the group law
and the formula for adding an integer to a floor.
\end{proof}


\begin{lemma}[Existence of admissible Liouville parameters]\label{checked-beta}
\lean{NilBohr.exists_admissibleBeta_between,NilBohr.AdmissibleBeta.liouville}
\leanok
Every nonempty interval inside $(0,1)$ contains an admissible
$\beta$, and each admissible $\beta$ is a Liouville number.
\end{lemma}
\begin{proof}
\proves{checked-beta}\leanok
Choose prime denominators using Bertrand's postulate at the
prescribed scale and round the numerators upward.
The rational approximants increase, while the approximants
plus the reciprocal of their denominators decrease. Their
supremum satisfies the one-sided approximation bounds.
The resulting error is smaller than every fixed inverse power
of the preceding denominator for sufficiently large indices.
\end{proof}

\begin{lemma}[Localized $U^S$ equality]\label{checked-cube}
\lean{NilBohr.cubeDifference_heisenberg_of_short,NilBohr.gowersCubePower_localized_vertical,NilBohr.normalizedGowers_localized_vertical,NilBohr.normalizedGowers_mul_polynomial_phase}
\leanok
For $S\ge3$, the exact equality (45) holds between the norm of
$f\chi$ with and without the localized vertical exponential,
using the finite zero-extended cube-power model. Multiplication
by an ordinary polynomial phase of degree less than $S$ also leaves
the normalized $U^S$ quantity unchanged.
\end{lemma}
\begin{proof}
\proves{checked-cube}\leanok
The finite cube calculation proves the phase factor cancels in every
nonzero term, and the same cube power and normalizer yield equality.
The multiplier proof of Lemma~4.5 combines this equality with the
Fourier expansion and the normalized Gowers triangle inequality.
For an ordinary polynomial, each additive difference lowers the
degree, so its $S$-fold difference vanishes and the exponential cube
factor is one.
\end{proof}

\begin{lemma}[Horizontal Diophantine and floor identities]\label{checked-floor}
\lean{NilBohr.horizontal_nonzero_alpha_bound,NilBohr.horizontal_zero_alpha_bound,NilBohr.bracketPhase_add_heisenbergPhase,NilBohr.floor_mul_eq_int_of_small_error,NilBohr.horizontal_scale_bounds,NilBohr.heisenberg_horizontal}
\leanok
The finite separation inequalities have a main term
$c_\alpha/(|k|vr^2)$ when $k\ne0$ and $1/r$ when $k=0$,
minus the corresponding approximation error $|l|v|\beta-a/r|$.
The bracket carry identity and small positive approximation
floor identity hold under the displayed arithmetic hypotheses.
\end{lemma}
\begin{proof}
\proves{checked-floor}\leanok
The finite separation uses nearest-integer distance and integer
divisibility. The bracket identity follows by splitting at
$\{\beta n\}=1-\beta$. The floor identity follows from the
unit interval above an integer. A second step checks the scale
uniformity for any approximating pair bracketing $X^{1/100}$.
\end{proof}

\begin{lemma}[Scalar Heisenberg equidistribution]\label{checked-scalar}
\lean{NilBohr.scalar_heisenberg_equidistribution,NilBohr.scalar_heisenberg_log_saving}
\leanok
The scalar estimate (25) follows from the fixed-Heisenberg
horizontal-character obstruction.
\end{lemma}
\begin{proof}
\proves{checked-scalar}\uses{checked-floor,external-inputs}\leanok
The finite-dilation lower bound rules out horizontal characters
of polylogarithmic size, even with prefix length $K\ge X^{1/3}$.
The hypothesis supplies the specialized Green--Tao criterion,
including its chart-mean and metric adaptation. The criterion
gives scalar equidistribution under the stated range.
\end{proof}

\begin{lemma}[Count splitting and density step]\label{checked-density}
\lean{NilBohr.selector_count_split,NilBohr.exists_density_increase_of_count_discrepancy,NilBohr.exists_densest_residue_progression,NilBohr.lift_scaled_configuration,NilBohr.density_iteration_contradiction,NilBohr.recurrence_of_density_increment_at_scale,NilBohr.heisenbergRecurrence_of_density_increment,NilBohr.density_increment_atom_from_count_inputs,NilBohr.density_model_count_lower,NilBohr.regularized_count_discrepancy,NilBohr.card_coincidentShiftParameters_le,NilBohr.norm_selected_count_le_exception_mass,NilBohr.log_initialModulus_eq_mangoldt_sum,NilBohr.exists_log_initialModulus_le_linear,NilBohr.eventually_log_initialModulus_at_density_le,NilBohr.exists_log_iteration_modulus_le_linear_in_w}
\leanok
The selector count splits by linearity. Given equal-mass factor
averages with sufficient count discrepancy, one of them
exceeds the prescribed density ceiling. There is a densest residue
progression, and rescaled configurations preserve the selected
prime. The logarithm of the initial modulus is the Mangoldt sum
through $w$, hence $O(w)$. A finite sequence of admissible
fixed-factor density increments contradicts density at most one.
\end{lemma}
\begin{proof}
\proves{checked-density}\uses{checked-rescaling}\leanok
Finite sums give the two-term weight split and product telescoping.
The mass identity gives the $\ell^1$ bound used in (20).
The affine identities lift every configuration point and the prime.
The prime powers dividing $\mathcal L(w)$ are exactly those at most
$w$. The primorial bound and the higher-prime-power estimate give
Chebyshev's $O(w)$ bound, including the initial finite cutoffs.
For the initial density and for every admissible stage this supplies
the linear-in-$w$ modulus budget used in the manuscript.
Logarithmic counting of the number of stages gives the finite
contradiction. The complete finite recursion takes Proposition~3.1
and a numerical parameter condition as premises. The later
asymptotic argument discharges that numerical condition. The
uniform density increment is also checked from Theorem~3.2
and KMTT weak regularity.
\end{proof}

\begin{lemma}[Local phases on one parameter cell]\label{checked-local-phase}
\lean{NilBohr.product_local_linear_phases,NilBohr.norm_local_phase_product_average_le,NilBohr.norm_partition_delete_le,NilBohr.norm_countingForm_le_local_tests_after_replacement}
\leanok
For fixed spatial $x$, the product of the local linear phases
on a single retained progression in $y$ is a polynomial phase
of degree at most $d$. Its average against the test weight is
bounded by the local polynomial-phase norm.
\end{lemma}
\begin{proof}
\proves{checked-local-phase}\leanok
The exponent is the sum of scalar multiples of the shifts and
constant terms. Polynomial degree bounds and finite products
give the stated reduction after the spatial blocks and residues
have been fixed. The finite selection and deletion bounds partition
$[M]$ into intervals of length $H$ and a tail, then count spatial
pairs within the polynomial displacement of a block boundary.
The retained count is bounded by the supremum of the local tests.
\end{proof}

\begin{lemma}[Sieve midpoint on a progression]\label{checked-sieve}
\lean{NilBohr.sieveMidpoint_bound,NilBohr.exists_modEq_of_dvd_affine,NilBohr.abs_card_dvd_affine_sub_le_one,NilBohr.primeProgression_as_affine,NilBohr.primeProgression_sieve_scales,NilBohr.cramerWeight_eq_sieveIndicator,NilBohr.cramerSieve_correlation_bound,NilBohr.primeDivisors_heisenberg_log_saving,NilBohr.eventually_cramerSieve_coefficients,NilBohr.MertensIntervalProductBound.eulerProduct_le_log}
\leanok
With upper and lower divisor envelopes, replacing the indicator by
their midpoint incurs at most half their average gap. For a
divisor coprime to the progression step, divisibility restricts to a
single residue class of parameter indices, whose count differs from
length divided by the divisor by at most one. The change to
$m=qy+1$ preserves the finite average and the degree of its
polynomial test. The resulting sieve progressions have the
uniform length and step bounds stated in Proposition~6.2. On the
residue-one progression the Cram\'er weight is its Euler product
normalization times the coprimality indicator omitting primes that
divide $q$. From the quoted sieve and interval-product assumptions,
the envelopes have absolute reciprocal-weighted gap at most
$(\log X)^{-D}$ for any fixed $D>0$ and all sufficiently large $X$.
\end{lemma}
\begin{proof}
\proves{checked-sieve}\uses{checked-twisted,external-inputs}\leanok
The midpoint inequality is pointwise followed by the triangle
inequality in a finite average. The arithmetic constraint is a
linear congruence with invertible coefficient.
The progression size estimates use the floor in $M$ and the
polylogarithmic bounds on $q$ and the first step. The finite
correlation bound sums the exceptional-divisor contribution using a
harmonic sum. On every other divisor progression that is coprime
to the first step, the twisted estimate applies to vertically
mean-zero tests. Deleting factors preserves the Mertens interval
bound. At level $\lfloor X^{1/20000}\rfloor$, the logarithmic
level ratio $U$ grows at least as a positive power of $\log X$.
The fundamental lemma then gives the stated envelope error.
\end{proof}

\begin{lemma}[Pointwise convergence and normalization steps]\label{checked-ergodic}
\lean{NilBohr.ae_mangoldtAverage_tendsto_zero_of_log_decay,NilBohr.ae_primeAverage_tendsto_zero_of_mangoldt_log_decay,NilBohr.ae_unrestrictedPrimeAverage_tendsto,NilBohr.OpenArc.exists_lipschitz_sandwich,NilBohr.centered_primeAverage_of_sandwich,NilBohr.ae_centered_selector_of_smooth_limits,NilBohr.pointwise_of_smooth_limits,NilBohr.pointwise_of_smooth_log_decay}
\leanok
For logarithmically bounded summands, a prescribed $L^1$ rate
$O((\log H)^{-6})$ implies almost-everywhere convergence to zero
of the prefix averages. Given the scalar counting and logarithmic
mass asymptotics, bounded log-weighted averages and unweighted
averages have the same limit. Scalar gap bounds turn centered lower
approximants into a centered interval weight. A positive prime-density
ratio and a centered comparison identify the selected prime limit.
\end{lemma}
\begin{proof}
\proves{checked-ergodic}\leanok
The first proof uses the manuscript's
$\lfloor\exp((j+1)^{1/3})\rfloor$ sampling and a gap bound.
Prime-power removal uses an explicit count and Mangoldt mass bound.
Log-weight removal uses the mass deficit of the weights $\log p/\log H$
conditional on the two scalar PNT asymptotics. The unrestricted
prime limit then follows conditionally on the weighted
Mirek--Wan--Wright statement. For interval weights the checked
Lipschitz ramps come from distance to the arc centre on the circle,
including arcs that cross the coordinate seam.
The last two deductions squeeze the centered interval average and
divide by the positive limiting count ratio. The full deduction
in Theorem~1.4 from the smooth logarithmic estimate is therefore
checked. The analytic bound at each positive exponent is checked from
Proposition~6.2 and the fixed-family estimate, and the bracket
deduction uses the rotation and cutoff lemmas below.
\end{proof}

\begin{lemma}[Orbit transference and duality]\label{checked-orbit}
\lean{NilBohr.orbit_counting_eq_average,NilBohr.integral_orbit_countingForm,NilBohr.integral_norm_weightedErgodicAverage_le_of_orbit_bound,NilBohr.orbit_bound_of_teravainen,NilBohr.exists_polynomial_orbit_window,NilBohr.mangoldt_L1_log_decay_of_primeLocalPhase}
\leanok
For integrable weighted averages and a common spatial window
containing all shifted orbit arguments, a uniform bound on
the discrete shifted count transfers to the $L^1$ norm of
the weighted ergodic average.
\end{lemma}
\begin{proof}
\proves{checked-orbit}\leanok
Average the measure-preserving change of variables over each
spatial index. Test against $\overline{A(x)}/|A(x)|$ on the
nonzero set of the weighted average $A$ and use zero otherwise.
This isolates the $L^1$ norm. For the fixed-family application the finite count estimate is the
explicit Ter\'av\'ainen Theorem~4.1 input. The common spatial window
uses a sum of absolute polynomial coefficients to bound all shifts.
The rate is supplied by Proposition~6.2 at $q=1$.
\end{proof}

\begin{lemma}[Spatial differencing and exceptional weights]\label{checked-difference}
\lean{NilBohr.countingForm_sq_le_spatial_energy,NilBohr.spatial_slice_square_eq_differences,NilBohr.exclude_exceptional_model_zero,NilBohr.weighted_local_inverse_of_kmtt_degree_uniform}
\leanok
Cauchy--Schwarz removes a bounded parameter weight, and grouping
spatial pairs by their difference gives the exact squared slice
in Lemma~5.1. The separate real-zero threshold comparison uses
the explicit Banks premise and the conductor bound.
\end{lemma}
\begin{proof}
\proves{checked-difference}\uses{external-inputs}\leanok
Use finite complex Cauchy--Schwarz and collect pairs by their
integer difference. For the real-zero step, multiply
$\log r\leq\log\widetilde Q$ by $1-\beta_0>0$.
The ambient prime model is an explicit KMTT input. The degree-uniform
KMTT local inverse hypothesis supplies a common bound before $C$ and
the family length are fixed. Spatial differencing and subsequent
phase extraction preserve that bound.
\end{proof}

\begin{lemma}[Product rotations and the bracket identity]\label{checked-rotations}
\lean{NilBohr.exists_quadratic_completion,NilBohr.ergodicSequence_circleRotationExtension,NilBohr.ae_quadratic_twisted_interval_of_smooth,NilBohr.ae_tendsto_primeAverage_fourier_quadratic_twist,NilBohr.phase_bracket_eq_piecewise,NilBohr.ae_bracket_fourier_cancellation,NilBohr.pointwise_bracket_of_heisenberg_smooth}
\leanok
The enlarged family and the circle-product system recover the
quadratic twist in (54). Equations (53) and (55) follow from
the smooth comparison on the product system. The bracket
carry decomposition preceding (56) also holds exactly.
\end{lemma}
\begin{proof}
\proves{checked-rotations}\uses{checked-ergodic,checked-floor}\leanok
Adjoin the absent degrees and extend the functions by constants.
The product system is probability preserving. Its character
is the desired quadratic twist times a circle factor of modulus
one. Fubini removes that factor. The floor formula gives both
branches of the bracket exponential. Fourier approximation and
endpoint ramps transfer these pieces to prime averages using
Vinogradov's scalar input. This gives (56) and the bracket
pointwise comparison.
\end{proof}


\begin{lemma}[Product obstruction and twisted prefixes]\label{checked-twisted}
\lean{NilBohr.choose_modEq_of_modEq_factorial_mul,NilBohr.realBinomialPhase_exp_totalVariation_le,NilBohr.heisenbergCoverTest_chartMean_eq_zero_of_lipschitz,NilBohr.norm_heisenbergTwistedAverage_le_of_product_obstruction,NilBohr.twisted_heisenberg_log_saving_of_dilated,NilBohr.twisted_heisenberg_log_saving_for_denominators,NilBohr.twisted_heisenberg_phaseNorm}
\leanok
Given the fixed scalar and product character criteria and explicit
bounds on the finite prefix scales, the polynomially twisted
Heisenberg average is bounded in terms of the two obstruction
thresholds and the variation of the residual twist. Under the
dilated scalar and horizontal estimates this yields a uniform
logarithmic saving at every exponent.
\end{lemma}
\begin{proof}
\proves{checked-twisted}\uses{checked-floor,checked-scalar,external-inputs}\leanok
The product criterion isolates a nonzero circle frequency.
Rounding in the integer binomial basis and splitting modulo
$|m|d!$ yields the covering-map expression and a slowly varying
error factor. Haar integration proves zero mean on the vertical
fibres of the covering test. The scalar criterion for each
sufficiently long prefix and partial summation give the finite bound.
The logarithmic scales then absorb the cover's distortion and the
discarded short prefixes. This deduction is for vertically mean-zero
tests on the quotient.
\end{proof}


\begin{lemma}[Selector cutoffs and discrepancy]\label{checked-selector-cutoffs}
\lean{NilBohr.quantitativeSelectorCutoffs,NilBohr.horizontalCircleEstimate_of_scalar,NilBohr.norm_countingForm_localized_selector_le,NilBohr.norm_countingForm_selector_boundary_le,NilBohr.norm_countingForm_selector_discrepancy_le}
\leanok
Three nonnegative short-arc cutoffs and a boundary cutoff give a
partition of unity. The Lipschitz bounds are $O(\varepsilon^{-1})$
and the boundary mean is $O(\varepsilon)$. The finite count
discrepancy splits into the localized prime-versus-model counts
and the boundary count.
\end{lemma}
\begin{proof}
\proves{checked-selector-cutoffs}\uses{checked-scalar,checked-cube,external-inputs}\leanok
Normalize the tent functions by their sum with a positive floor.
Outside the endpoint strip the floor has no effect, and the
Lipschitz and measure bounds follow from the lower bound on the
denominator and the length of a ball. Horizontal projection of the
scalar orbit estimate bounds the mean on the finite progression.
The localized multiplier and normalized Gowers laws then control
the prime discrepancy, using the Bernstein Fourier bound.
\end{proof}


\begin{lemma}[Centered chart test and seam removal]\label{checked-centered}
\lean{NilBohr.heisenbergCenteredTest_orbit,NilBohr.heisenbergCenteredTest_verticalMean_zero,NilBohr.heisenbergCenteredTest_lipschitz,NilBohr.heisenberg_seamAverage_of_horizontal,NilBohr.centeredHeisenbergPhaseEstimate_of_obstructions}
\leanok
A Lipschitz circle weight centered at its Haar mean yields a
right lattice-invariant, vertically mean-zero test after cutting off
the horizontal seam. Its quotient Lipschitz norm is
$O_F(1+\varepsilon^{-1})$. The scalar and twisted orbit estimates
then give Equation~(27) uniformly.
\end{lemma}
\begin{proof}
\proves{checked-centered}\uses{checked-floor,checked-scalar,checked-twisted,external-inputs}\leanok
The orbit formula uses the vertical chart coordinate
$z-x\lfloor y\rfloor$. A relative coordinate bound for each group
increment controls the change in the centered test. If the horizontal
coordinate crosses the chart seam, the cutoff bounds the surviving
values by the horizontal displacement divided by $\varepsilon$.
Summing over chains and taking the quotient infimum gives the
Lipschitz estimate. Circle translation gives the vertical mean, and a
horizontal majorant bounds the discarded seam average.
\end{proof}

\begin{lemma}[Prime-weighted seam error]\label{checked-prime-seam}
\lean{NilBohr.primeProgression_mangoldt_seam_error}\leanok
At the seam width $\eta=(\log X)^{-A-2}$, replacing the centered
selector on the shifted-prime progression by the full cutoff chart
test costs $O_F((\log X)^{-A})$ against $\Lambda$, uniformly in the
polynomial phase.
\end{lemma}
\begin{proof}
\proves{checked-prime-seam}\uses{checked-centered,checked-sieve}\leanok
At divisor one the prime progression has length at least $X^{1/2}$
and step at most $X^{1/10000}$. The common horizontal denominator
is coprime to its step. The scalar estimate bounds the seam mass
by $O(\eta)$. The von Mangoldt weight costs at most $\log X$, and the
polynomial phase has modulus one.
\end{proof}

\begin{lemma}[Cram\'er cancellation with the full chart test]\label{checked-prime-cramer}
\lean{NilBohr.norm_cramerPhaseAverage_le_of_sieve,NilBohr.eventually_siftedPhase_error,NilBohr.primeCramerPhase_log_saving}\leanok
The full centered chart test at seam width $(\log X)^{-E}$,
for any fixed $E>0$, has arbitrary logarithmic saving on the
prime progression against the Cram\'er model, uniformly in the
polynomial test of degree at most $d$.
\end{lemma}
\begin{proof}
\proves{checked-prime-cramer}\uses{checked-sieve,checked-centered}\leanok
Use the exact affine reindexing and normalize the chart test by a
fixed sup bound. On every good divisor progression, all the chart's
vertical fibres have mean zero and its quotient Lipschitz norm
is polylogarithmic. Harmonic summation, the exceptional divisors,
and the sieve gap give one more inverse logarithm than the target,
before multiplication by the Euler product.
\end{proof}

\begin{lemma}[Fourier restriction on a progression]\label{checked-fourier-restriction}
\lean{NilBohr.localized_prime_model_replacement_progression,NilBohr.residue_indicator_fourier_l1,NilBohr.intervalCyclicFourierCoefficient_l1_le,NilBohr.quantitativePrimeSelectorCutoffs,NilBohr.localized_prime_model_replacement_full_chart,NilBohr.prime_model_replacement_centered_chart,NilBohr.primeModelReplacement_log_saving_of_kmtt}
\leanok
A localized and polynomially twisted progression average of an
ambient weight is controlled by its ambient $U^S$ quantity, with the
length ratio $X/|P'|$, the Fourier mass of the circle cutoff, and
a factor bounded by $2+\log X$. The bound extends to the full
centered seam test and yields arbitrary logarithmic saving for the
ambient prime-model discrepancy on the progressions of Proposition~6.2.
\end{lemma}
\begin{proof}
\proves{checked-fourier-restriction}\uses{checked-cube,checked-selector-cutoffs,external-inputs}\leanok
Use cyclic Fourier orthogonality to write the two interval and
congruence indicators as linear phases. The residue coefficients
have mass one. Geometric summation bounds each nonzero interval
coefficient by the reciprocal of twice its centered frequency,
so the total mass is at most one plus the harmonic sum to $X$.
The Gowers triangle and mean bounds and the checked phase
equalities give the normalized restriction.
The bound is uniform in the vertical Fourier index. Absolute Fourier
convergence and an exact short-arc partition of the seam cutoff give a
polylogarithmic loss for the full chart test. The prime-model bound
absorbs this loss.
\end{proof}

\begin{lemma}[Scalar and product Heisenberg models]\label{checked-group-model}
\uses{selector}
\lean{NilBohr.heisenbergLeibmanBasis_rational,NilBohr.heisenbergProductLeibmanBasis_rational,NilBohr.heisenbergNilmanifoldMetric_topology,NilBohr.heisenbergProductNilmanifoldMetric_topology,NilBohr.malcevQuotientChainDistance_eq_heisenbergQuotientDistance,NilBohr.malcevQuotientChainDistance_eq_heisenbergProductQuotientDistance,NilBohr.isLeibmanHaar_heisenbergNilmanifoldMeasure,NilBohr.isLeibmanHaar_heisenbergProductNilmanifoldMeasure,NilBohr.heisenbergNilmanifoldMeasure_unique,NilBohr.heisenbergProductNilmanifoldMeasure_unique,NilBohr.heisenbergTwist_productLipschitzBound}\leanok
The scalar and product quotients have adapted Mal'cev bases of
rationality height at most one. Their chain metrics induce the
quotient topologies. Unit-cube volume pushes forward to the normalized
Haar probabilities and gives the chart integral on each. The
dilated scalar orbit is filtered polynomial of degree at most two.
Adjoining an ordinary polynomial of degree at most $d$ gives the
product of filtration degree $\max(2,d)$. The twist
$\Phi(g)\ee(-u)$ is lattice invariant, has mean zero after circle
integration, and increases the scalar supremum-plus-Lipschitz
bound by at most an absolute factor.
\end{lemma}
\begin{proof}
\proves{checked-group-model}\uses{checked-floor}\leanok
The global matrix charts are Euclidean, with polynomial group laws.
The lattices are closed and discrete and the closed unit cubes
map onto the compact Hausdorff quotients. Smooth one-parameter axes
give second-kind coordinates $(x,y,z-xy)$ and the extra real coordinate.
The sole nonzero scalar Lie bracket of basis axes is
$[X,Y]=Z$ up to sign from exchanging the axes. Ordering the extra axis
before $Z$ when $d=1$ and after $Z$ for $d\ge2$ makes the basis
adapted to the appropriate product filtration.

Small chain displacements control the coordinate differences and
conversely. The quotient infimum over the lattice gives the actual
quotient metric because the lattice is closed. For Haar, take the
pushforward of normalized volume on the half-open fundamental cube.
Shears preserve ambient Lebesgue measure, and periodicity gives the
chart formula. Uniqueness follows by lifting an invariant quotient
probability with counting measure on the lattice fibres and comparing
the resulting ambient Haar measures. The scalar and circle
projections are bounded by the product chain distance, giving the
product test bound by expanding the difference of two products.
\end{proof}

\begin{lemma}[Quantitative Leibman specialization]\label{checked-green-tao-adapter}
\lean{NilBohr.scalarHeisenbergObstruction_of_greenTao,NilBohr.productHeisenbergObstruction_of_greenTao,NilBohr.SelectorLiteratureInputs.of_quantitativeLeibman,NilBohr.PrimePhaseLiteratureInputs.of_quantitativeLeibman}\leanok
The general quantitative Leibman input gives the scalar and product
character criteria used in Proposition~4.3 on the respective
nilmanifolds of Lemma~\ref{checked-group-model}.
The recurrence and prime-phase input bundles can therefore be
assembled directly from the listed source interfaces.
\end{lemma}
\begin{proof}
\proves{checked-green-tao-adapter}\uses{external-inputs,checked-group-model}\leanok
Use the descended scalar test on the dilated orbit and the test
$\Phi(g)\ee(-u)$ on the product orbit. Their means and Lipschitz bounds
are those of the concrete model. In the product application a fixed
factor in the test bound is absorbed by lowering the tolerance.
A horizontal character vanishes on the commutator axis. Its
remaining integer coordinates give the horizontal linear form
and, in the product case, the coefficient of the real polynomial.
First differences in the scalar case and the iterated differences
at zero in the product case give the required binomial-coefficient
inequalities. Constants are chosen after the dimension and
filtration degree and before the orbit and test parameters.
\end{proof}

\begin{lemma}[Finite identities for Appendix A]\label{checked-separating-finite}
\uses{selector,d:Cq}
\lean{NilBohr.phaseNorm_dilatedWeight,NilBohr.parameterLength_separatingFamily,NilBohr.countingForm_separatingFamily_dilate,NilBohr.countingForm_separatingFamily_repeat,NilBohr.residueCorrelation_le_phaseNorm,NilBohr.abs_heisenbergPhase_sub_separatingQuadratic_le}\leanok
Zero extension to $q\Z$ gives a phase norm of exactly $q^{-1}$
times the original phase norm. The fixed-family count normalization
contributes $q^2$, or $q$ when the spatial inputs repeat on all
$q$ residue classes. For a rational approximation with
$0\le\beta-a/r\le E$ and $E qM<1$, the error in the quadratic
approximation to $Z(qy)$ on $r\mid y$ is at most
$|\alpha|E(qM)^2/2$.
\end{lemma}
\begin{proof}
\proves{checked-separating-finite}\uses{checked-floor,checked-fourier-restriction}\leanok
Affine changes of variable transport the real polynomial tests in
both directions. The residue expansion has Fourier mass one, so a
correlation on the multiples of $r$ contributes its value divided
by $r$ to the global phase-norm lower bound. On those multiples
the small nonnegative approximation error fixes the floor exactly.
The count identities follow by reindexing both summation variables
and tracking the numerator and denominator separately.
\end{proof}

\begin{lemma}[The separating scales]\label{checked-separating-scales}
\uses{selector,d:Cq}
\lean{NilBohr.SeparatingScale,NilBohr.eventually_separatingScale,NilBohr.tendsto_separatingLengthQuotient,NilBohr.tendsto_separatingLength,NilBohr.eventually_separatingStep,NilBohr.eventually_separatingRelativeModulus,NilBohr.eventually_separatingRelativeScaleWindow,NilBohr.tendsto_separatingCoefficient_div_lengthQuotient,NilBohr.tendsto_separatingApproximation,NilBohr.eventually_separatingLogModulus_logLogLength,NilBohr.tendsto_logDenominator_div_logSeparatingModulus,NilBohr.tendsto_separatingLogLength_div_exp_sqrt_log,NilBohr.tendsto_logSeparatingScaleX_div_logNext,NilBohr.eventually_initialModulus_dvd_separatingModulus,NilBohr.tendsto_separatingW_div_logSpatialScale_pow,NilBohr.tendsto_logSeparatingModulus_div_sqrt_logSpatialScale}\leanok
The floored power of $\mathcal L(w)$ and the integer multiple
of the denominator in the definition of $M$ give positive
scales, $\log q\asymp(\log r_j)^2$, and
$q\le M$. The logarithmic ratios give
$\log M\sim\tfrac14\exp(\sqrt{\log r_j})$ and
$qM=r_{j+1}^{1/4+o(1)}$. At $X=qM$, the denominator bracket and the bounded
powers of the progression step are valid for large $j$.
\end{lemma}
\begin{proof}
\proves{checked-separating-scales}\uses{checked-separating-finite}\leanok
The floor changes $\log q$ by at most $w\log w$. The next
denominator grows faster than any fixed power of the current
one or of the chosen modulus. The resulting polynomial bounds
in $r_{j+1}$ control $H$, $M$ and $X$ on the same sequence.
\end{proof}

\begin{lemma}[Local tests at the separating scales]\label{checked-separating-local}
\uses{selector,d:Cq}
\lean{NilBohr.SeparatingLocalPhaseDecay,NilBohr.separatingLocalPhase_log_bound,NilBohr.eventually_separating_test_length,NilBohr.separatingCount_log_bound,NilBohr.separatingGlobalPhase_tendsto_zero,NilBohr.separatingLocalPhase_of_obstructions,NilBohr.separatingCount_of_kmtt}\leanok
The centered progression estimate gives arbitrary logarithmic
saving in the local phase tests at the separating scales.
Every progression required by Theorem~4.2 at $\rho=(\log M)^{-1}$
meets the local length threshold. Testing $[M]$ in particular
gives convergence to zero for its ordinary global $u^3$ norm.
\end{lemma}
\begin{proof}
\proves{checked-separating-local}\uses{p:heisenberg-phase,t:main,checked-separating-scales}\leanok
Reparametrize the integer progressions and transport the
polynomial tests along the affine change. The exact cardinality
bound for the residue class loses at most one. The small powers
of $q$ leave the required length margin.
\end{proof}

\begin{lemma}[Continuous circle tests]\label{checked-separating-fourier}
\lean{NilBohr.exists_nonzero_fourierCoeff_of_nonconstant,NilBohr.circleMean_centered_mul_fourier,NilBohr.tendsto_circle_average_of_fourier}\leanok
A nonconstant continuous circle function has a nonzero Fourier
coefficient at some nonzero frequency. The mean of its centered
function times the matching conjugate character is that coefficient.
Cancellation of every nonzero mode on finite rows gives
equidistribution under continuous tests.
\end{lemma}
\begin{proof}
\proves{checked-separating-fourier}\leanok
Use uniqueness of circle Fourier coefficients for the first
claim and trigonometric polynomial approximation for the last.
\end{proof}

\begin{lemma}[Quadratic correlation on the denominator multiples]\label{checked-separating-quadratic}
\uses{selector,d:Cq}
\lean{NilBohr.prefixAverage_tendsto_zero_of_shift_correlations,NilBohr.separatingQuadratic_modes_tendsto_zero,NilBohr.separatingQuadratic_mode_sq_le_growing_shift,NilBohr.tendsto_separatingCoefficient_div_lengthQuotient_pow,NilBohr.separatingGlobalPhase_lower_bound}\leanok
At the separating scales the phase on the denominator multiples
is approximated by a quadratic. Its nonzero Fourier modes have
vanishing averages. The correlation with a quadratic test and
the residue expansion give a positive multiple of $r_j^{-1}$
as a lower bound for the full-interval phase norm.
\end{lemma}
\begin{proof}
\proves{checked-separating-quadratic}\uses{checked-separating-finite,checked-separating-scales,checked-separating-fourier}\leanok
The geometric-series and bad-approximation bounds control each
nonzero shift in the quadratic sequence. Growing-shift van der
Corput gives the displayed $H^{-1/3}+|k|a_jq^2r_jH^{-2/3}$ bound,
whose terms vanish at the separating scales. A second checked
proof takes fixed shifts first. Trigonometric approximation gives
the limit, and the surviving constant is a nonzero Fourier coefficient.
\end{proof}

\begin{lemma}[A full Gowers lower bound]\label{checked-separating-gowers}
\uses{selector,d:Cq}
\lean{NilBohr.separating_gowers_lower_bound}\leanok
For every fixed $S\ge3$, scalar Heisenberg equidistribution on
the long, coprime progression and the localized vertical
multiplier give a positive lower bound for $U^S[M]$.
\end{lemma}
\begin{proof}
\proves{checked-separating-gowers}\uses{p:heisenberg-phase,l:vertical-mode-multiplier,checked-separating-scales,checked-separating-fourier}\leanok
The product of the horizontal cutoff mean and a nonzero
vertical Fourier coefficient is the limit of the tested
average. The Gowers mean bound and Lemma~4.5 bound the
average by a constant times the full Gowers norm.
\end{proof}
